\phantomsection
\chapter*{Conclusion}\label{chap:Conclusion}
\addcontentsline{toc}{chapter}{Conclusion}

In the corpus study we found a general preference for syntactic agreement (except for profession nouns). We also found significant effects of target-controller ordering in the whole dataset, an effect of distance for socially gendered nouns and an effect of type of agreement target for profession nouns. The corpus study led to conducting two acceptability judgement experiments: one with split hybrids (depreciative forms) and the other with lexical hybrids (profession nouns). There we found an effect of the type of noun, meaning that items with depreciative forms had significantly lower acceptability ratings than items with profession nouns. There was, however, no interaction with semantic agreement, so there was no definite conclusion regarding the preference for semantic agreement with lexical or split hybrids. In the experiment using profession nouns we also found that items with relative pronouns were significantly more acceptable than items with predicate agreement, but here again there was no interaction with semantic agreement. The studies show that semantic agreement in acceptable in Polish and sensitive to several factors of Corbett’s Agreement Hierarchy.

The results of the studies show directions for further research: experimental studies with more participants and more types of agreement targets should be conducted to accurately verify the hypotheses regarding hybrid gender agreement. Additionally, agreement with depreciative forms should be studied more, as it is an example of split hybrids, which are a part of a newer extension of the Agreement Hierarchy and are not very common. Studying those forms can give important insight into the rules of hybrid agreement, but it should be done carefully, as these forms are colloquial and speakers may be reluctant to accept them in experimental settings.
